%% ma: L12-B08-0170
%% lop: 12
%% chuong: 2
%% bai: 8
%% dang: 12.8.53
%% loai: TN4
%% muc: VD
%% dap_an: "A"
%% nguon: "CĐ2.4.pdf giáo viên gửi 27/10/2026 (M07, Lê Minh Khôi – Nguyễn Anh Khôi, Chương II Chủ đề 4 Các bài toán về tâm tỉ cự), Câu 20"
%% kien_thuc: "Bài 8 (khoảng cách hai điểm, biểu thức tọa độ), Bài 6 (bình phương vô hướng, tâm tỉ cự)"
%% trang_thai: chinh-thuc
%% ghi_chu: "Giáo viên đã duyệt 27/10/2026. Đề gốc dùng chung giả thiết cho Câu 20–22; đã chép vào từng câu"
\begin{ex}
Trong không gian $Oxyz$, cho ba điểm $M(2;1;-2)$, $N(-1;0;3)$, $P(1;-3;5)$ và điểm $I$ bất kỳ trong không gian. Tìm giá trị nhỏ nhất của biểu thức $K=IM^2+IN^2+4IP^2$.
\choice{\True $K_{\min}=\dfrac{367}6$}{$K_{\min}=\dfrac{367}2$}{$K_{\min}=\dfrac{367}4$}{$K_{\min}=\dfrac{367}{16}$}
\loigiai{Gọi $E$ thỏa $\overrightarrow{EM}+\overrightarrow{EN}+4\overrightarrow{EP}=\vec 0\Rightarrow E=\dfrac{M+N+4P}6=\left(\dfrac56;-\dfrac{11}6;\dfrac72\right)$.

$K=6IE^2+EM^2+EN^2+4EP^2\ge EM^2+EN^2+4EP^2=\dfrac{367}6$.}
\end{ex}
